% ============================================================================
%  Tema de prueba en escala de grises. Reemplaza cada color de marca por su
%  luminancia (Rec. 709), que es aproximadamente lo que entrega una impresora
%  monocroma. Se activa con la opción «grises» de la clase:
%      \documentclass[grises]{audit-oferta}
%  Sirve para revisar que nada dependa del color para leerse. No se entrega.
% ============================================================================

\definecolor{audit-marino}{gray}{0.16}
\definecolor{audit-turquesa}{gray}{0.62}
\definecolor{audit-gris-claro}{gray}{0.92}
\definecolor{audit-tinta}{gray}{0.14}
\definecolor{audit-secundario}{gray}{0.40}
\definecolor{audit-filete}{gray}{0.81}
\definecolor{audit-fila}{gray}{0.96}

\definecolor{fig-nube}{gray}{0.30}
\definecolor{fig-sitio}{gray}{0.40}
\definecolor{fig-terminal}{gray}{0.36}
\definecolor{fig-flota}{gray}{0.32}
\definecolor{fig-red}{gray}{0.28}

% Todo color definido con xcolor sale en grises. Las imágenes y los diagramas
% importados conservan su color, porque son archivos externos.
\selectcolormodel{gray}
